\documentclass[12pt]{article}
\usepackage[utf8]{inputenc}
\usepackage[T1]{fontenc}
\usepackage{hyperref}
\title{Bokha Language Resources for K-12 Curriculum Development}
\author{Compiled by Claude Sonnet 5.5 (claude-sonnet-5-5) for Andr\'as Kornai}
\date{}
\begin{document}
\maketitle

\section{Introduction}
Bokha (ISO 639-3 ybk; also Akapa, Aphu, Black Muji, Bokho, Flowery Phula) is a Sino-Tibetan language of China. Wikitongues lists it as threatened (EGIDS 6b). Joshua Project places the Bokha within China's official Yi nationality, says their language is not mutually intelligible with the Yi varieties around them, and estimates about 12{,}000 speakers. Joshua Project's Muji, Northern entry (ymx) also lists ``Bokha'' as an alternate name, so sources may confuse the two. My searches found very little educational material.

\section{Primary Educational Resources}
No primers, textbooks or school materials in Bokha turned up in my searches. Joshua Project lists the written status of the language as unknown.

\section{Linguistic and Reference Materials}
No dictionary or grammar turned up in my searches. Language-profile pages only:
\begin{itemize}
\item Wikitongues language page for Bokha: \url{https://wikitongues.org/languages/ybk/}
\item Joshua Project language page (ybk): \url{https://joshuaproject.net/index.php/languages/ybk}
\end{itemize}

\section{Religious and Community Resources}
Joshua Project reports that the Bokha still need a Bible translation.
\begin{itemize}
\item Joshua Project, Bokha in China: \url{https://joshuaproject.net/people_groups/18417}
\end{itemize}

\section{Digital and Audio Resources}
No audio, video or corpus resources turned up in my searches.

\section{Curriculum Development Context}
There is no evidence of an orthography, of school use or of published teaching materials. A curriculum would first need fieldwork and orthography development with the community, and the community's own priorities would have to guide it. Teaching would probably take place alongside Mandarin.
\end{document}